\scanpage{3}
Carte sphérique \inserted{(orientée)} avec repère (dans cat. \unsure{isotopiques|isotopique})

\[
\rotatebox{90}{$\simeq$}
\]
morphisme $\mathbb P^{1}_{\mathbb C}\overset{f}{\unsure{\longrightarrow|-}}\mathbb P^{1}_{\mathrm C}$ ``spécial'' i.e
\[
 f(0)=0,\quad f(1)=1,\quad f(\infty)=\infty.\quad f\ \text{étale au dessus de }\mathbb P^{1}-\{0,1,\infty\}
\]
\[
\rotatebox{90}{$\simeq$}
\]
fonction rationnelle $f(z)$ telle que $f(0)=\deleted{\unsure{1}}0$, $f(1)=1$, $f(\infty)=\infty$ et telle que pour $z\in\mathbb C$ \deleted{telle que}\inserted{avec} $f(z)\deleted{\in}\notin\{0,\infty\}$ $f'(z)=0$, on ait $f(z)\deleted{\in}=1$ \deleted{\unsure{$\{0,1,\infty$}}$\}$ \inserted{et de plus} $f''(z)\ne0$ (i.e. la ramification de $f$ en $z$ est $2$)\ldots.

\[
\rotatebox{90}{$\simeq$}
\]
couple $P,Q$ \inserted{$[f=P/Q\quad !\,]$} de polynômes $\in\mathbb C(z)$ premiers entre eux \inserted{(i.e. sans zéros communs)} tels que $P(0)=0$, (donc $Q(0)\ne0$) $P(1)=Q(1)$ (donc $\ne0$), $\deg Q<\deg P$ (ces conditions expriment resp\textsuperscript{t} $f(0)=0$, $f(1)=1$, $f(\infty)=\infty$)
% slanted
\margin{On pourrait aussi normaliser \unsure{au lieu de cel|au lieu de cela} par $Q(0)=1$, ou $P$ unitaire, ou $Q$ unitaire\ldots}
normalisés (par exemple) par $P(1)=Q(1)=1$, et tels que pour $z\in\mathbb C$ avec $[f(z)\deleted{\in}\notin\{0,\infty\}, f'(z)=0,\ \mathrm{i.e.}\,]$ \inserted{$P(z)\ne0$,} $Q(z)\ne0$,
\[
 \frac{P'(z)\deleted{Q(z)}}{P(z)}-\frac{Q'(z)}{Q(z)}=0,
\]
on ait $P(z)=Q(z)$, et de plus
\[
 \frac{P''(z)}{P(z)}-\frac{Q''(z)}{Q(z)}\ne0\qquad\text{i.e.}\quad P''(z)\ne Q''(z)\ldots
\]
[ NB on a
\[
\begin{aligned}
 Q(z)P'(z)-P(z)Q'(z)&=\deleted{Q(z)[P'(z)-P(z)]+P(z)(Q(z)-Q'(z))}\\
 &\phantom{=}\inserted{(Q(z)-P(z))P'(z)+P(z)(P'(z)-Q'(z))}
\end{aligned}
\]
donc si \inserted{pour $P(z);Q(z)\ne0$,} cette quantité est nulle, (i.e. $f'(z)=0$), alors $P(z)=Q(z)$ \deleted{\unsure{au|ou}} i.e $Q(z)-P(z)=0$ implique $P'(z)=Q'(z)$ et inversement si $P'(z)=Q'(z)$ on a $(Q(z)-P(z))P'(z)=0$ donc \uline{si} $P'(z)\ne0$, on a $P(z)=Q(z)$ ! Mais peut-on avoir $P'(z)=0$ (d'où aussi $Q'(z)=0$) \inserted{avec $P'(z)=Q'(z)$}\unsure{$\}$|)} sans avoir $P(z)=Q(z)$ ?

\deleted{\unsure{Réfléchissons} $Q(T)=$}

Prenons p.ex. $Q(T)=1$, les conditions sur $P$ deviennent $P(0)=0$, $P(1)=1$, $\deg P\geq1$ \inserted{(conséquence des deux précédentes)}, donc $P(T)=TP_0(T)$ où on a $P_0(1)=1$, la question devient si on peut avoir \inserted{(pour $P(z)\ne0$)} $P'(z)=0$ sans avoir $P(z)=1$ certes oui p.ex. $\deleted{P(T)=T^2}$ si $P(T)=T(aT+b)=aT^2+bT$, $a+b=1$, on a
% margin note written sideways
\margin{\raggedright $P'(T)=2aT+b$ qui est nul \inserted{(supposons $a\ne0$)} pour $z=-\frac{b}{2a}$ \deleted{sera} $\ne\left\{\begin{smallmatrix}0\\-\frac{b}{a}\end{smallmatrix}\right.$ \deleted{\illegible{1}} les deux zéros de $P$ si $b\ne0$, et $P(-\frac{b}{2a})=-\frac{b}{2a}(-\frac b2+b)=-\frac{b^2}{4a}$ qui \deleted{\unsure{peut être}} est $\ne1$ \deleted{\unsure{en prenant}} si $b^2\ne-4a$\unsure{---}.}
% archive sheet number: 2
